\section{VideoSQL}

\subsection{Index Building}
We used the same index sturcture as MAGIC-Video \cite{li_bridging_2026}, which includes video captions, entities, visual clips, event chain, etc. Readers can referenced to MAGIC-Video for more details. However, the original pipeline is complext and involving many slow stages like OpenIE. We modified the indexing process such that an VLM (Qwen 3.8) to directly output these information based on the input.

\subsection{Agentic SQL Queries}
MAGIC-Video uses Personalized Page Rank (PPR) to retrieve nodes in the graph by a query generated by the LLM, with which the agent has no control over the relations between nodes. We still keep a simple semantic search tool for the model to search all text and images using embeddings. For example, if the agent wants to find who ate the apple yesterday at noon, it likely issues a query like "person ate apple yesterday noon" and hopes the PPR returns relevant nodes. Instead, we formulate this using a relational database. Given the same index structure as MAGIC-Video, the agent can write a SQL query like in \cref{lst:sql-example}. Even though searches like this could be replaced with dedicated tools like \texttt{search\_semantic(object="apple", time\_s=time\_s, time\_end=time\_end)}, using SQL as an interface provides a unified interface with more freedom and does not require specific design for each case (for example, if we need to search "the person who ate the apple that is not Luke", we would need to re-design the tool signature). It can also be seen as a form of Programmatic Tool Calling. Such an interface can also handle complex queries using advanced SQL features like JOIN, GROUP BY, and sub-queries when needed.

Our work differs from SAG \cite{wu_sag_2026} in that the LLM controls the SQL generation and it can use all SQL features, unlike SAG which only used JOIN as a way to expand the seed nodes.

\begin{lstlisting}[language=SQL, caption={An example SQL query for finding who ate the apple within a given time window.}, label={lst:sql-example}]
SELECT *
FROM semantic_triples
WHERE predicate IN ('eat', 'eats', 'ate', 'eating')
  AND object IN ('apple', 'the apple', 'an apple')
  AND snapshot_s >= $start_s
  AND snapshot_s < $end_s;
\end{lstlisting}


\subsection{Trainable Skills}

